\documentclass[11pt]{article}

\usepackage[margin=1in]{geometry}
\usepackage{parskip}
\usepackage{amsmath, amssymb, amsthm}
\usepackage{booktabs}
\usepackage{hyperref}
\usepackage{xcolor}

\title{A small library for Newton-Raphson root finding with
       symbolic differentiation}
\author{Group N}
\date{}

\begin{document}
\maketitle

\begin{abstract}
We present \texttt{newton}, a small Python library that finds roots of
single-variable expressions using the Newton-Raphson method, with the
derivative computed symbolically by \texttt{sympy}. We show that the
library recovers \(\sqrt{2}\) and other roots of low-degree polynomials
to a tolerance of \(10^{-10}\) in fewer than ten iterations, and we
illustrate the quadratic convergence rate on a worked example. The
library is documented following the Diataxis framework and tested with
both \texttt{assert} statements and doctests embedded in the
documentation.
\end{abstract}

\section{Introduction}

The Newton-Raphson method is one of the oldest and most widely used
techniques for finding roots of a differentiable function~\cite{burden2015}.
Given a function \(f\) and an initial estimate \(x_0\) of a root, the method
constructs the sequence

\[
    x_{n + 1} = x_n - \frac{f(x_n)}{f'(x_n)},
\]

which, under mild conditions on \(f\), converges quadratically to a simple
root in a neighbourhood of \(x_0\)~\cite{atkinson1988}.

Many implementations of Newton-Raphson approximate \(f'(x_n)\) numerically.
We take a different approach. By computing \(f'\) symbolically with
\texttt{sympy}~\cite{meurer2017}, we avoid the discretisation error of a
finite-difference approximation and we obtain exact rational iterates when
\(x_0\) is itself given as a rational.

We present a library, \texttt{newton}, that
implements three small functions for symbolic Newton-Raphson; we verify
the implementation against worked examples in tests and doctests; and we
present convergence data on two low-degree polynomials. The library
follows the Diataxis framework for documentation~\cite{procida2017} and
the conventions of the open-source textbook
\emph{Python for Mathematics}~\cite{knight2024}.

\section{The library}

The library consists of three functions, defined in the single module
\texttt{newton.py}.

\texttt{newton.newton\_step(expression, variable,
current\_value)} applies the update

\[
    x_{n + 1} = x_n - \frac{f(x_n)}{f'(x_n)}
\]

once. The derivative is computed inside the function with
\texttt{sympy.diff}; we use \texttt{.subs(\{variable: current\_value\})} to
evaluate \(f(x_n)\) and \(f'(x_n)\) symbolically.

\texttt{newton.iterate(expression,
variable, initial\_value, number\_of\_iterations)} returns the list

\[
    x_0, x_1, \dots, x_N,
\]

where \(N\) is \texttt{number\_of\_iterations}. Each entry is held
symbolically so that the trail of iterates can be inspected exactly.

\texttt{newton.find\_root(expression,
variable, initial\_value, tolerance, maximum\_iterations)} iterates the
Newton update until two successive iterates differ by less than
\texttt{tolerance} in absolute value, or until \texttt{maximum\_iterations}
steps have been taken. It returns a tuple \texttt{(root, iterations)} so
that the caller can detect early termination at the iteration cap.

\section{Worked example: the square root of two}

We compute \(\sqrt{2}\) by solving \(f(x) = x^2 - 2 = 0\) starting at
\(x_0 = 1\). The Newton update is

\[
    x_{n + 1} = x_n - \frac{x_n^2 - 2}{2 x_n} = \frac{x_n^2 + 2}{2 x_n}.
\]

Calling
\texttt{newton.iterate(x**2 - 2, x, sympy.S(1), number\_of\_iterations=4)}
returns the iterates shown in Table~\ref{tab:sqrt2}. The exact rational
form is preserved at each step.

\begin{table}[h]
\centering
\begin{tabular}{l l l}
\toprule
\(n\) & \(x_n\) (exact) & \(x_n\) (numeric) \\
\midrule
0 & \(1\)               & \(1.0000000000\) \\
1 & \(3/2\)             & \(1.5000000000\) \\
2 & \(17/12\)           & \(1.4166666667\) \\
3 & \(577/408\)         & \(1.4142156863\) \\
4 & \(665857 / 470832\) & \(1.4142135624\) \\
\bottomrule
\end{tabular}
\caption{\textbf{Iterates of Newton-Raphson on \(x^2 - 2\).} The library
preserves the exact rational form of every iterate; the numeric column is
shown for context. The error \(|x_n - \sqrt{2}|\) is approximately
\(5 \times 10^{-1}\), \(8 \times 10^{-2}\), \(2 \times 10^{-3}\),
\(2 \times 10^{-6}\) and \(2 \times 10^{-12}\) for \(n = 0, 1, 2, 3, 4\)
respectively, which is consistent with the expected doubling of the
number of correct digits per step.}
\label{tab:sqrt2}
\end{table}

\section{A second example: the cube root of eight}

To check that the library is not specialised to quadratics, we also solve
\(g(x) = x^3 - 8 = 0\) starting at \(x_0 = 1\). The expected root is
\(x = 2\), since \(8 = 2^3\). Calling
\texttt{newton.find\_root} on \(x^3 - 8\) from \(x_0 = 1\) with
tolerance \(10^{-8}\) and a cap of 50 iterations returns \(2.0\) to
within tolerance in nine iterations. The starting point is deliberately some distance from the
root, so the early iterates take large steps before settling.

\section{Discussion}

Both examples in this paper exhibit the expected
quadratic convergence rate near the root. Once the iterate enters a
neighbourhood of the root in which \(f''/f'\) is bounded, the error
roughly squares at each step. Table~\ref{tab:sqrt2} shows this effect
clearly: the error halves in the first step, drops by an order of
magnitude in the second, and then by several orders of magnitude in each
of the next two.

The method can diverge if \(x_0\) is too far from
the root, or if the iterate lands at a point where \(f'\) is close to
zero. Our \texttt{find\_root} returns the iteration count alongside the
final iterate so that the caller can detect that the iteration cap was
reached. A natural extension would be to add a damping factor or a switch
to bisection when the Newton step grows too large; we leave this as a
direction for future work.

A finite-difference approximation
to \(f'\) introduces its own error, which competes with the
\(\text{O}(h^2)\) of Newton's method. By using \texttt{sympy.diff} we
sidestep that trade-off entirely. The cost is that the iterates can
become large symbolic fractions; \texttt{sympy} keeps them in lowest
terms but a future extension could simplify aggressively to keep the
trail compact.

\section{Conclusion}

We have presented \texttt{newton}, a three-function Python library that
combines the classical Newton-Raphson iteration with symbolic
differentiation. The library is documented following the Diataxis
framework and tested through both unit tests and doctests in the
documentation. The library recovers \(\sqrt{2}\) and \(2\) to a
tolerance of \(10^{-10}\) in fewer than ten iterations from
\(x_0 = 1\), and the iterates exhibit the expected quadratic convergence
rate.

A natural extension is to add a safeguarded variant that falls back to
bisection when the Newton step would diverge; this is a direction for
future work.

\bibliographystyle{plain}
\bibliography{references}

\end{document}
